\documentclass[11pt,a4paper]{article}
\usepackage[margin=0.85in]{geometry}
\usepackage{amsmath,amssymb}
\usepackage{booktabs}
\usepackage{array}
\usepackage{enumitem}
\usepackage{hyperref}
\usepackage{microtype}

\hypersetup{
    colorlinks=true,
    linkcolor=black,
    citecolor=blue,
    urlcolor=blue
}

\begin{document}

% =========================================================================
% 1. TITLE OF THE PROJECT & AUTHOR DETAILS
% =========================================================================
\title{\textbf{\Large ModelBench: A Web-Based Platform for Empirical Machine Learning Model Evaluation and Comparative Performance Benchmarking}}

\author{
    \textbf{Ujjwal Singh} (Roll No: 262IS036) \qquad \textbf{Aryan} (Roll No: 262IS007) \\[0.4em]
    \normalsize M.Tech I Year (CSE-IS), Semester I \\[0.2em]
    \normalsize Course: CS801 Computing Lab $\mid$ Course Instructor: Prof. P. Santhi Thilagam \\[0.2em]
    \normalsize Department of Computer Science and Engineering \\
    \normalsize National Institute of Technology Karnataka, Surathkal
}
\date{\textbf{Submission Deadline:} September 04, 2026}

\maketitle
\thispagestyle{empty}

\vspace{-1.2em}

% =========================================================================
% 2. DESCRIPTION
% =========================================================================
\section{Description}
\label{sec:description}
Machine Learning (ML) algorithms are increasingly deployed to solve practical analytical and predictive problems across diverse fields. However, selecting an appropriate model for a given task requires testing different algorithms, tuning their parameters, and evaluating their behavior on common data.

\textbf{ModelBench} is an interactive, web-based software platform designed for systematic training, evaluation, and side-by-side comparison of classical machine learning algorithms. Through an accessible interface, users can choose curated real-world datasets, apply standard preprocessing techniques, configure model hyperparameters, and directly compare the performance of multiple algorithms under identical data conditions. The platform produces clear statistical metrics, diagnostic visualizations, and an interactive prediction tool to help users understand the strengths and limitations of different models.

% =========================================================================
% 3. MOTIVATION
% =========================================================================
\section{Motivation}
\label{sec:motivation}
When developing machine learning solutions for real-world problems---such as medical risk assessment, property price estimation, or customer churn prediction---practitioners rarely rely on a single algorithm. Multiple candidate models must be evaluated to identify the best balance of accuracy, generalization, and computational efficiency.

In typical practice, these experiments are conducted across separate Jupyter notebooks or standalone scripts. This ad-hoc approach presents several practical challenges:
\begin{enumerate}[leftmargin=*,itemsep=1pt,topsep=2pt]
    \item \textbf{Inconsistent Preprocessing and Data Splits:} Running models in isolated scripts often leads to unintentional variations in feature scaling, data cleaning, or train/test splits, making fair comparison difficult.
    \item \textbf{Over-Reliance on Single Summary Metrics:} Experiments often focus only on overall accuracy or mean loss, overlooking critical error distributions, confusion matrices, or precision-recall trade-offs.
    \item \textbf{Lack of Controlled Comparison Tools:} Manually collecting, recording, and charting results across multiple models is repetitive, time-consuming, and prone to human error.
\end{enumerate}

\textbf{ModelBench} provides a centralized and consistent testing environment where multiple ML algorithms can be trained, evaluated, and compared systematically on the same benchmark datasets.

% =========================================================================
% 4. PROBLEM STATEMENT AND OBJECTIVES
% =========================================================================
\section{Problem Statement and Objectives}
\label{sec:problem_statement}

\subsection{Problem Statement}
To design and implement a web-based, modular software platform that provides standardized data preprocessing, model training, performance evaluation, and comparative benchmarking of classical machine learning algorithms on curated real-world datasets.

\subsection{Objectives}
The specific objectives of the project are:
\begin{enumerate}[leftmargin=*,itemsep=1pt,topsep=2pt]
    \item \textbf{Automated Preprocessing:} Implement standard data validation, missing value handling, feature scaling, and reproducible 80/20 train/test splitting.
    \item \textbf{Support for Core ML Algorithms:} Integrate representative algorithms across three major learning categories:
    \begin{itemize}[leftmargin=*,itemsep=0pt,topsep=1pt]
        \item \textit{Regression:} Linear Regression, Multiple Linear Regression, Polynomial Regression, Ridge Regression.
        \item \textit{Classification:} Logistic Regression, $K$-Nearest Neighbors (KNN), Decision Trees, Random Forests, Support Vector Machines (SVM).
        \item \textit{Unsupervised Learning:} $K$-Means Clustering, Principal Component Analysis (PCA).
    \end{itemize}
    \item \textbf{Standardized Evaluation:} Compute appropriate statistical metrics for each task ($R^2$, RMSE, MAE for regression; Accuracy, Precision, Recall, F1-score, ROC-AUC for classification; Silhouette score for clustering).
    \item \textbf{Diagnostic Visualizations:} Generate clear graphical charts, including residual error plots, confusion matrices, ROC curves, and decision boundaries.
    \item \textbf{Side-by-Side Model Comparison:} Provide comparative summary tables and charts to contrast multiple models trained on the same dataset.
    \item \textbf{Interactive Prediction (What-If Simulation):} Provide simple slider-based inputs to test live model predictions interactively.
\end{enumerate}

% =========================================================================
% 5. SOLUTION DESIGN (TENTATIVE)
% =========================================================================
\section{Solution Design (Tentative)}
\label{sec:solution_design}
The proposed platform follows a modular, client-server design organized into five sequential workflow steps:

\begin{enumerate}[leftmargin=*,itemsep=2pt,topsep=2pt]
    \item \textbf{Dataset Selection:} The user chooses a curated real-world dataset representing a specific domain:
    \begin{itemize}[leftmargin=*,itemsep=0pt,topsep=1pt]
        \item \textit{Healthcare:} Diabetes or heart disease prediction (Classification).
        \item \textit{Real Estate / Economics:} Housing price or medical insurance cost estimation (Regression).
        \item \textit{Customer Analytics:} Customer segmentation or churn analysis (Clustering / Classification).
    \end{itemize}
    \item \textbf{Data Preprocessing:} The backend automatically verifies feature types, handles missing values, scales numerical features, and applies reproducible data splitting.
    \item \textbf{Model Training:} The user selects an algorithm and adjusts its key hyperparameters (such as regularization strength, tree depth, or number of neighbors). The backend trains the model using established Python scientific libraries.
    \item \textbf{Evaluation and Diagnostic Charting:} The system calculates standard performance metrics and generates diagnostic plots, which are returned as interactive web charts.
    \item \textbf{Comparative Analysis and Live Testing:} Users can review multiple models side-by-side on a common scoreboard or use dynamic input sliders to test live predictions.
\end{enumerate}

% =========================================================================
% 6. DIFFERENT PHASES OF THE PROJECT
% =========================================================================
\section{Different Phases of the Project}
\label{sec:phases}
The project will be executed in four structured phases aligned with the course timeline:

\begin{itemize}[leftmargin=*,itemsep=2pt,topsep=2pt]
    \item \textbf{Phase 1: Problem Definition \& Baseline Setup (Weeks 1--3):} Finalize project scope, curate benchmark datasets, and build baseline regression and classification models in Jupyter notebooks.
    \item \textbf{Phase 2: Core Supervised Models \& MidSem Demonstration (Weeks 4--6):} Implement Multiple Linear, Polynomial, Ridge, Logistic Regression, and KNN models. Connect the ML engine to a lightweight web interface for the MidSem 40\% live code demonstration.
    \item \textbf{Phase 3: Decision Trees, SVM \& Unsupervised Models (Weeks 7--9):} Add Decision Trees, Random Forests, Support Vector Machines, $K$-Means clustering, and side-by-side model comparison tools for Continuous Evaluation 2.
    \item \textbf{Phase 4: What-If Simulator, System Refinement \& Final Report (Weeks 10--11):} Implement the interactive prediction simulator, perform end-to-end testing, and prepare the comprehensive $\LaTeX$ final project report.
\end{itemize}

% =========================================================================
% 7. IMPLEMENTATION LEVEL DETAILS
% =========================================================================
\section{Implementation Level Details}
\label{sec:implementation}

\begin{table}[!ht]
\centering
\small
\caption{System Implementation Technologies and Tools}
\label{tab:tech_stack}
\begin{tabular}{@{}llp{8.2cm}@{}}
\toprule
\textbf{Layer / Component} & \textbf{Technology} & \textbf{Function in Project} \\ \midrule
Frontend Interface & React.js, Tailwind CSS & User interface, parameter inputs, and responsive layout. \\
Data Visualization & Plotly.js & Rendering interactive charts, confusion matrices, and residual plots. \\
Backend Application & FastAPI (Python 3.10+) & Request routing, input validation, and API management. \\
Machine Learning Engine & Scikit-learn, NumPy, Pandas & Data preprocessing, model training, and metric calculations. \\
Development \& Reporting & \LaTeX, Git, GitHub & Technical report compilation, version control, and code tracking. \\ \bottomrule
\end{tabular}
\end{table}

% =========================================================================
% 8. EXPECTED OUTCOMES
% =========================================================================
\section{Expected Outcomes}
\label{sec:outcomes}
Upon successful completion of the project, the expected outcomes are:
\begin{enumerate}[leftmargin=*,itemsep=1pt,topsep=2pt]
    \item A functional, responsive web application for training, evaluating, and comparing classical ML algorithms.
    \item A standardized collection of curated benchmark datasets with documented feature sets.
    \item Implementation of 8+ classical machine learning models spanning regression, classification, and clustering.
    \item A comparative benchmarking dashboard displaying side-by-side metrics and diagnostic graphs.
    \item An interactive What-If feature testing tool for real-time model prediction.
    \item Complete technical documentation and $\LaTeX$ reports detailing the experimental methodology and findings.
\end{enumerate}

% =========================================================================
% 9. TENTATIVE TIMELINES
% =========================================================================
\section{Tentative Timelines}
\label{sec:timelines}
The project milestones are aligned with the official NITK CS801 Computing Lab evaluation schedule:

\begin{table}[!ht]
\centering
\small
\caption{Project Milestones and Evaluation Schedule}
\label{tab:timelines}
\begin{tabular}{@{}llp{8.2cm}@{}}
\toprule
\textbf{Milestone / Phase} & \textbf{Target Date} & \textbf{Planned Deliverables} \\ \midrule
Extended Abstract Submission & September 04, 2026 & Finalized 3-page $\LaTeX$ abstract submission and review. \\
Continuous Evaluation 1 & September 21, 2026 & 10-minute presentation on problem analysis and tentative design. \\
\textbf{MidSem Evaluation} & \textbf{October 12, 2026} & $\LaTeX$ progress report, presentation, and \textbf{40\% live code demonstration}. \\
Continuous Evaluation 2 & November 01, 2026 & Progress presentation and \textbf{50\% live code demonstration}. \\
\textbf{EndSem Evaluation} & \textbf{Nov 10--13, 2026} & Comprehensive $\LaTeX$ final report and \textbf{60\% complete code demonstration}. \\ \bottomrule
\end{tabular}
\end{table}

% =========================================================================
% 10. REFERENCES
% =========================================================================
\section{References}
\label{sec:references}
\begingroup
\renewcommand{\section}[2]{}
\begin{thebibliography}{9}
\small
\bibitem{pedregosa2011scikit}
F.~Pedregosa \textit{et al.}, ``Scikit-learn: Machine learning in Python,'' \textit{Journal of Machine Learning Research}, vol.~12, pp.~2825--2830, 2011.

\bibitem{geron2019hands}
A.~G{\'e}ron, \textit{Hands-On Machine Learning with Scikit-Learn, Keras, and TensorFlow: Concepts, Tools, and Techniques to Build Intelligent Systems}, 2nd~ed. Sebastopol, CA: O'Reilly Media, 2019.

\bibitem{tiangolo2018fastapi}
S.~Ram{\'\i}rez, ``FastAPI: Modern, high-performance web framework for building APIs with Python,'' 2018. [Online]. Available: \url{https://fastapi.tiangolo.com/}

\bibitem{uci_repository}
D.~Dua and C.~Graff, ``UCI Machine Learning Repository,'' School of Information and Computer Sciences, University of California, Irvine, 2017. [Online]. Available: \url{http://archive.ics.uci.edu/ml}
\end{thebibliography}
\endgroup

\end{document}
